\section{Consequences, comparisons, and limitations}
\label{sec:consequences}

\subsection{A new structural target}

Earlier conditional routes in the project parameterize exact computation by
frontier width, twist-run count, extremal Khovanov windows, transfer rank, or
hierarchy data.  The present result adds a parameter on the diagram-rewriting
side:
\[
  (\text{birth number},\ \text{unlock depth}).
\]
It is independent of homological width.  A diagram may have a difficult exact
scan but a small causal RIII certificate, or conversely a simple exact invariant
calculation but no short monotone simplification.

This separation suggests an adaptive policy: estimate local RIII branch growth
and exact-scan frontier growth, then spend a bounded amount on whichever has the
more favorable certified parameter.  Such a policy must preserve the existing
rule that abandoned work cannot affect verdicts.

\subsection{Why the result is stronger than a depth-only bound}

A naive global depth-$k$ RIII search has $O(N)$ choices at each level and hence
$N^{O(k)}$ behavior.  At $k=O(\log N)$ this is already
$N^{O(\log N)}$, quasi-polynomial, but with a large exponent and no recognition
of local causal structure.  More importantly, the current implementation does
not perform that complete global search because it would be too expensive.

The causal theorem decomposes the exponent.  Only births pay a factor of $N$;
all connected continuations pay $O(k)$.  When $b$ is constant and
$k=O(\log N)$, the exponent improves from $O(\log N)$ to
$O(\log\log N)$.  This is the main asymptotic gain.

\subsection{Why polynomial Reidemeister length is not enough}

Lackenby's polynomial move bound permits crossing increases and controls a much
broader class of sequences.  A polynomial total length $N^{11}$ inserted into
our enumeration formula would be useless.  The present theorem needs two finer
facts:
\begin{enumerate}
\item the neutral run between decreases is short;
\item only a few support-disjoint branches are born during that run.
\end{enumerate}
Proving these facts for all unknots would be a major additional theorem and may
be false for the strictly crossing-nonincreasing move language.

\subsection{Potential relation to planar separators}

The active RIII footprint contains $O(k)$ darts and induces a small subgraph of
the projection.  If this active region admits balanced separators that are
stable under local moves, one could attempt a divide-and-conquer search in the
style of frontier algorithms.  Such a theorem might replace the $N^b$ birth
factor by a smaller separator-dependent term.  At present, however, a birth can
occur anywhere, and no safe separator theorem for the dynamic RIII causal graph
is proved here.

\subsection{Potential relation to hierarchy certificates}

A short RIII causal trace is a two-dimensional diagram certificate; an essential
hierarchy is a three-dimensional complement certificate.  There may be a
transfer principle: a uniformly local hierarchy simplification could constrain
where diagrammatic births occur, while a small causal footprint might identify
a bounded part of a handle structure that needs modification.  The 2026
hierarchy algorithm makes this a timely direction, but no such bridge is
claimed.

\subsection{Lower-bound questions}

The birth parameter also supports negative results.  To prove that this route
cannot solve the unrestricted problem, one could construct unknot diagrams for
which every crossing-nonincreasing first-unlocking trace has either large depth
or large birth number.  The connectedness theorem says a shortest trace cannot
contain causally irrelevant branches, so a lower bound on births would be a
meaningful topological obstruction rather than a search artifact.

\subsection{Soundness boundary}

The preprocessor has one-sided semantics:
\[
  \text{successful replayable trace}\Rightarrow\Unk,
  \qquad
  \text{failure or timeout}\Rightarrow\Unknown.
\]
It cannot certify knottedness.  Knottedness remains the responsibility of sound
invariants or exact complement/homology methods.  This boundary should be
preserved in every CLI and serialized-certificate interface.

\subsection{Assessment of the research contribution}

Within its scope, the report delivers four concrete advances.
\begin{enumerate}
\item It replaces an empirical locality heuristic by an explicit two-parameter
      completeness theorem.
\item It proves a restricted quasi-polynomial and polynomial classification.
\item It identifies and repairs a subtle but natural false proof strategy.
\item It supplies code and a low-risk opt-in integration path.
\end{enumerate}
The unresolved step is geometric: establish useful upper bounds on births and
depth for broad, natural families of unknot diagrams, or design a different
certificate when those parameters are large.
